%% ---------------------------------------------------------------------------
%% LQL-Equiv -- supplementary material for the SoftwareX submission
%% Compiled separately from main.tex and uploaded as a single PDF.
%% Compile with:  pdflatex supplementary ; bibtex supplementary ;
%%                pdflatex supplementary ; pdflatex supplementary
%% ---------------------------------------------------------------------------
\documentclass[preprint,12pt,number]{elsarticle}

\usepackage[T1]{fontenc}
\usepackage[utf8]{inputenc}
\usepackage[a4paper,top=2.4cm,bottom=2.4cm,left=2.4cm,right=2.4cm]{geometry}
\usepackage{amsmath,amssymb}
\usepackage{booktabs}
\usepackage{graphicx}
\usepackage{float}
\renewcommand{\topfraction}{0.9}
\renewcommand{\textfraction}{0.1}
\usepackage{url}
\usepackage[hidelinks]{hyperref}

\journal{SoftwareX}

%% Sections are numbered S1, S2, ... and figures S1, S2, ...
\renewcommand{\thesection}{S\arabic{section}}
\renewcommand{\thefigure}{S\arabic{figure}}
\renewcommand{\theequation}{S\arabic{equation}}

\begin{document}

\begin{frontmatter}
\title{Supplementary material for\\
       \emph{LQL-Equiv: a browser-based implementation of
       linear-quadratic-linear dose equivalence for radiotherapy}}
\author[label1]{Cyril Voyant}
\author[label2]{Daniel Julian}
\affiliation[label1]{organization={Mines Paris, PSL University},
            city={Sophia-Antipolis}, country={France}}
\affiliation[label2]{organization={Centre de Canc\'erologie du Grand Montpellier},
            city={Montpellier}, country={France}}
\end{frontmatter}

This document holds two analyses referred to from the main text. Section~S1
states what the 2014 release computed and what version 3.0 computes, and measures
the difference. Section~S2 compares the software against published randomised
trials and dose-tolerance reviews. Both are reproduced in the repository as
\texttt{docs/COMPARISON-2014.md} and \texttt{docs/CLINICAL-PLAUSIBILITY.md}, with
the scripts that generate every number.
\section{What changed between the 2014 release and version 3.0}
\label{app:2014}

Version 3.0 does not reproduce the 2014 application for the organ at risk. This
appendix states what each version computes, in its own notation, and measures the
difference. The tumour calculation is unchanged in substance.

Everything below can be checked against the original, which is openly archived at
\url{https://github.com/cyrilvoyant/LQ-Equiv} and at
\url{https://doi.org/10.5281/zenodo.16739883}. The block discussed here is
\texttt{pushbutton4\_Callback} in \texttt{eqd\_matlb.m}, inside
\texttt{codesource.zip}, at lines 3186--3245. The captured outputs of the running
application ship as \texttt{tests/data/golden.jsonl}, so the comparison can be
repeated without MATLAB.

\subsection*{Notation, repeated here so that this document stands alone}

A course delivers $n$ fractions of dose $d$ over an overall time $T$, and $d_r$
is the reference fraction size, conventionally 2\,Gy. Write
\[
L(x) \;=\; x\left[1 + \frac{(1+H_m)\,x}{\alpha/\beta}\right]
\quad (x < d_t),
\qquad
L(x) \;=\; d_t\left[1 + \frac{d_t}{\alpha/\beta}\right]
  + \frac{\gamma}{\alpha}\,(x - d_t)
\quad (x \geq d_t),
\]
for the per-fraction term below and above the transition dose
$d_t = 2\,\alpha/\beta$, where $H_m$ is the incomplete-repair correction, zero
unless two fractions are given a day. Treatment runs five days a week, so $s$
treatment sessions span
\begin{equation}
\Theta(s) \;=\;
\begin{cases}
s, & s < 6,\\[0.6ex]
s + 2\left(\left\lfloor \dfrac{s-6}{5} \right\rfloor + 1\right), & s \geq 6,
\end{cases}
\label{eq:calendar}
\end{equation}
calendar days, weekends included. Proliferation is charged per elapsed day,
under the two models the 2014 paper gives the two tissues:
\begin{equation}
P_\mathrm{target}(T) = D_\mathrm{prol}\,(T - T_k)_{+},
\qquad
P_\mathrm{organ}(T) = D_\mathrm{prol}\,T ,
\qquad
D_\mathrm{prol} = \frac{\ln 2}{\alpha\,T_p},
\label{eq:loss}
\end{equation}
Dale with a kick-off time for the target, Van Dyk \cite{vandyk1989} without one
for the organ. The biologically effective dose of a course spanning $[t_0, t_1]$
is then
\begin{equation}
\mathrm{BED}(d, n; t_0, t_1) \;=\; n\,L(d) - \big[P(t_1) - P(t_0)\big],
\label{eq:bed}
\end{equation}
and the equivalent dose is obtained by solving, for $n_r$,
\begin{equation}
\mathrm{BED}\big(d_r,\,n_r;\,\tau,\,\tau + \Theta(n_r)\big)
 \;=\; \mathrm{BED}\big(d,\,n;\,t_0,\,t_1\big),
\qquad
\mathrm{EQD}_{d_r} = n_r\,d_r ,
\label{eq:core}
\end{equation}
$\tau$ being the reference time already spent by preceding courses. Additivity
follows from $P$ telescoping across consecutive segments, whatever calendar is
used.

\subsection*{What the 2014 source computes}

On a grid of hundredths of a fraction, the routine first solves
\begin{equation}
n_r \;=\; \arg\min_{x} \left|\;
\Big[n L(d)   - \tfrac{7}{5}\,n\,D_\mathrm{prol}\Big]
-\Big[x L(d_r) - \tfrac{7}{5}\,x\,D_\mathrm{prol}\Big] \right| ,
\label{eq:legacy-search}
\end{equation}
so that proliferation is charged on both sides at the rate of $7/5$ days per
fraction. It then reports
\begin{equation}
\mathrm{EQD}^{\,2014}_{d_r} \;=\; n_r\,d_r
\;-\;\big[\Theta(n + g) - \Theta(n_r)\big]\,D_\mathrm{prol} .
\label{eq:legacy-report}
\end{equation}
The kick-off time $T_k$ appears nowhere in \eqref{eq:legacy-search}, and that is
correct: the 2014 paper gives the two tissues two different proliferation models
and says so. Its equations (3) and (4) apply Dale to the target volume, with a
kick-off time; its equations (6) and (7) apply Van Dyk \cite{vandyk1989} to the
organ at risk,
where a recovered dose accrues from the first day and, in the paper's words, ``the
kick-off time is no longer considered''. Version 3.0 keeps that distinction.

What \eqref{eq:legacy-report} adds has no counterpart in the paper. Equation (12)
there defines the equivalent dose as $2 n_0$, the fraction count times the
reference dose, with nothing after it, and the equivalent dose the same routine
reports for the tumour a few lines further down is exactly that.

\subsection*{What version 3.0 computes}

Equations \eqref{eq:calendar}--\eqref{eq:core} as printed above: the two
published proliferation models, one elapsed calendar shared by both sides of the
equality with weekends and missed sessions included, and
$\mathrm{EQD}_{d_r} = n_r d_r$ with nothing added afterwards.

\subsection*{The criterion}

A convention returns a different number that is consistent with its own
definition; an inconsistency returns a number that is not. The definition here is
\eqref{eq:core}, so the test is whether the reported dose, delivered at $d_r$,
reproduces the biologically effective dose it replaced. For a rectum
($\alpha/\beta = 3.9$\,Gy, $D_\mathrm{prol} = 0.3$\,Gy\,d$^{-1}$) receiving
20~$\times$~3\,Gy:

\begin{center}
\small
\begin{tabular}{@{}lrr@{}}
\toprule
 & \textbf{2014} & \textbf{version 3.0} \\
\midrule
BED of the schedule (Gy)                      & 97.75  & 97.75 \\
Equivalent dose reported (Gy)                 & 82.69  & 75.25 \\
\quad as fractions of 2\,Gy                   & 41.34  & 37.63 \\
BED of that schedule (Gy)                     & \textbf{107.74} & \textbf{97.75} \\
\quad against 97.75\,Gy                       & $+9.97$ & $0.00$ \\
\bottomrule
\end{tabular}
\end{center}

The 2014 value fails the test, and the reason is visible in
\eqref{eq:legacy-search}--\eqref{eq:legacy-report}: $n_r$ already balances the
proliferation of the two schedules, so subtracting the difference in their spans
a second time charges it twice. Of the 7.65\,Gy that
\eqref{eq:legacy-report} adds for this case, 7.36\,Gy is the amount already
priced into $n_r$ and 0.30\,Gy is the genuine refinement from replacing the rate
$7/5$ by the staircase. Over 7850 deliverable schedules the round trip closes for
28.6\,\% of the 2014 values and for all of version 3.0's.

\subsection*{One difference that is a convention}

The choice between the staircase $\Theta$ and its long-run rate $7n/5$ is not
settled by any such test: both are defensible, and they never differ by more than
2.4 days, worth 0.7\,Gy per side at $D_\mathrm{prol} = 0.3$\,Gy\,d$^{-1}$.
Version 3.0 uses the staircase throughout, because proliferation is charged per
elapsed day and weekends are elapsed days: 40 sessions delivered Monday to
Friday span 54 days, not the 56 the rate returns. The 2014 code used the rate
inside \eqref{eq:legacy-search} and the staircase in \eqref{eq:legacy-report},
which is the mixture that made the double charge hard to see. Using the rate
consistently instead would move the case above from 75.25\,Gy to 75.03\,Gy.

The staircase costs something in return. It drops the reference BED by two days
of proliferation at every weekend, so a target falling in that drop is met by two
fraction counts, one on each side of the jump, both exact. The shorter is
reported, so that the equivalent dose remains a function of its input.

\subsection*{How large the difference is}

Figure~\ref{fig:difference2014} takes 14\,712 deliverable plans: one, two and
three successive courses, 1 to 12\,Gy per fraction, once and twice a day, with
interruptions up to fourteen missed sessions in total, over all 34 organs and 20
target sites of the library, keeping only those whose equivalent dose falls
between 2 and 100\,Gy EQD2.

\begin{center}
\small
\begin{tabular}{@{}lrr@{}}
\toprule
Version 3.0 minus 2014       & \textbf{Organ at risk} & \textbf{Target} \\
\midrule
Plans                        & 12\,622 & 14\,543 \\
Unchanged ($|\Delta| < 0.01$\,Gy) & 27.5\,\% & 42.8\,\% \\
Median                       & $-0.67$\,Gy & $-0.00$\,Gy \\
5th, 95th percentile         & $-12.51$, $+4.63$\,Gy & $-1.41$, $+6.20$\,Gy \\
Range                        & $-44.3$ to $+35.2$\,Gy & $-3.8$ to $+36.7$\,Gy \\
\bottomrule
\end{tabular}
\end{center}

The organ moves and the target barely does, which is what
\eqref{eq:legacy-report} predicts: the second term exists only on the organ
branch, and what remains on the target is the calendar alone. The difference is
nil at the reference fraction size, where the two calendars coincide and
\eqref{eq:legacy-report} adds nothing, and nil for the eight organs whose
$D_\mathrm{prol}$ is zero; it grows with the departure from $d_r$. It is mostly
negative, the 2014 release having reported more dose to the organ than the
schedule carries, and turns positive for fraction sizes below $d_r$ and for
schedules prolonged by an interruption. Splitting a course into two moved the
2014 total by up to 0.6\,Gy; version 3.0 is additive by construction.

\begin{figure}[t]
\centering
\includegraphics[width=\textwidth]{figures/fig5_difference2014.pdf}
\caption{Equivalent dose, version 3.0 minus the 2014 release, over 14\,712
deliverable plans. (a) Organ at risk, on a logarithmic count axis, stacked by
number of successive courses. (b) Both tissues, split by whether the plan
contains an interruption and by fraction regime. The target moves only through
the calendar; the organ carries the second term as well.}
\label{fig:difference2014}
\end{figure}

A value published before 2026 can be recovered with
\texttt{Options.legacy\_2014()}, which reproduces
\eqref{eq:legacy-search}--\eqref{eq:legacy-report} together with the three other
other defects listed in the article. It is what the 4438-schedule
non-regression suite measures; nothing else in the software uses it.

%% ---------------------------------------------------------------------------
\section{External comparison against the clinical literature}
\label{app:trials}

The numerical checks in the article are internal: they establish that
the software solves its own equations. This appendix asks a different question.
Randomised trials have compared fractionations directly, so the schedules a trial
found non-inferior should come out close in equivalent dose, and those it did not
should not. Every value below was recomputed with the released version, reference
dose 2\,Gy per fraction, prostate $\alpha/\beta = 3.1$\,Gy and $d_t = 6.2$\,Gy,
spinal cord $\alpha/\beta = 2$\,Gy, $d_t = 4$\,Gy and $\gamma/\alpha = 5$.

This is a coherence check and not a clinical validation. Non-inferiority in a
trial depends on the population, the dosimetry, the volumes, the margins, the
imaging and the systemic treatment, none of which a fractionation model sees.

\subsection*{Target volume: prostate}

\begin{center}
\small
\begin{tabular}{@{}llrrl@{}}
\toprule
Trial & Schedule & EQD2 & Against & Trial outcome \\
\midrule
HYPO-RT-PC & 42.7\,Gy / 7  & 77.94 & 78.00 & non-inferior \\
CHHiP      & 60\,Gy / 20   & 72.46 & 74.00 & non-inferior \\
CHHiP      & 57\,Gy / 19   & 68.79 & 74.00 & not demonstrated \\
PROFIT     & 60\,Gy / 20   & 72.46 & 78.00 & non-inferior \\
PACE-B     & 36.25\,Gy / 5 & 73.34 & 76.12, 78.00 & non-inferior \\
\bottomrule
\end{tabular}
\end{center}

\paragraph{HYPO-RT-PC, agreement to within 0.06\,Gy.} The software places
42.7\,Gy in seven fractions at 77.94\,Gy EQD2 against 78.00\,Gy for the
conventional arm. The trial found 84\,\% failure-free survival in both arms at
five years, adjusted hazard ratio 1.002 (95\,\% CI 0.758--1.325)
\cite{widmark2019hypo}, and 72\,\% against 65\,\% at ten years, adjusted hazard
ratio 0.84 (95\,\% CI 0.69--1.03) \cite{nilsson2026hypo}. This is the closest
external agreement obtained, with one caveat: $6.1 < d_t = 6.2$\,Gy, so the case
exercises the quadratic branch and the time correction and not the linear tail.

\paragraph{CHHiP, the ordering recovered.} 60\,Gy in 20 fractions comes out
1.54\,Gy below the conventional arm, and the trial reported 90.6\,\% against
88.3\,\% failure-free at five years, hazard ratio 0.84 (90\,\% CI 0.68--1.03),
$p_{NI} = 0.0018$. 57\,Gy in 19 fractions comes out 5.21\,Gy below, and there the
trial reported 85.9\,\%, hazard ratio 1.20 (90\,\% CI 0.99--1.46), $p_{NI} =
0.48$: non-inferiority was not established \cite{dearnaley2016chhip}. The
software orders the two hypofractionated arms as the trial did. Most CHHiP
patients received androgen suppression, which limits what the comparison can
carry radiobiologically.

\paragraph{PROFIT, a quantitative disagreement.} The same 60\,Gy in 20 fractions
sits 5.54\,Gy below 78\,Gy in 39, yet PROFIT, without hormone therapy, reported
85\,\% biochemical or clinical failure-free survival in both arms, hazard ratio
0.96 (90\,\% CI 0.77--1.20), and demonstrated non-inferiority
\cite{catton2017profit}. This is not an arithmetic error but a sensitivity to
$\alpha/\beta$. Under the plain linear-quadratic form with
$\alpha/\beta = 1.5$\,Gy,
\[
\mathrm{EQD2} = 60\,\frac{3 + 1.5}{2 + 1.5} = 77.14\ \mathrm{Gy},
\]
within a gray of the conventional arm. A pooled clinical estimate of
$\alpha/\beta = 1.4$\,Gy (0.9--2.2) has been published for prostate
\cite{miralbell2012alphabeta}. The tabulated 3.1\,Gy cannot be treated as
universal.

\paragraph{PACE-B, the linear tail exercised.} At 7.25\,Gy per fraction the dose
is above $d_t$, so the linear tail applies:
\[
\mathrm{BED_{LQL}} = 5\left[6.2\left(1+\tfrac{6.2}{3.1}\right)
  + 5\,(7.25-6.2)\right] = 119.25\ \mathrm{Gy_{3.1}},
\qquad
\frac{119.25}{1 + 2/3.1} = 72.49\ \mathrm{Gy},
\]
and 73.34\,Gy once the reference schedule carries its own overall time. That
places the stereotactic arm 2.78\,Gy below 62\,Gy in 20 and 4.66\,Gy below
78\,Gy in 39, while the trial reported 95.8\,\% against 94.6\,\% failure-free at
five years, hazard ratio 0.73 (90\,\% CI 0.48--1.12), $p_{NI} = 0.004$
\cite{vanas2024paceb}. The linear tail is conservative for this schedule.

\subsection*{Organ at risk: spinal cord}

\paragraph{Stereotactic, 25.3\,Gy in five fractions.} HyTEC associates a maximum
dose of 25.3\,Gy in five fractions with a 1--5\,\% risk of radiation myelopathy
\cite{sahgal2021hytec}. At 5.06\,Gy per fraction the software gives
\[
\mathrm{BED_{LQL}} = 5\left[4\left(1+\tfrac{4}{2}\right) + 5\,(5.06-4)\right]
  = 86.50\ \mathrm{Gy_2},
\qquad
\mathrm{EQD2} = \frac{86.50}{1 + 2/2} = 43.25\ \mathrm{Gy},
\]
against 44.65\,Gy under the plain linear-quadratic form, a difference of
$-1.40$\,Gy or $-3.1$\,\%. The reported complication probability is 1.45\,\%,
numerically inside the HyTEC range. That agreement should not be pressed: HyTEC
defines its constraint on a maximum dose to the cord or thecal sac, and the
software has no volume information at all.

\paragraph{Conventional, 50\,Gy in 25 fractions.} The equivalent dose is exactly
50.00\,Gy, as it must be at the reference fraction size, but the reported
complication probability is 5.13\,\%. QUANTEC puts the risk of myelopathy near
0.2\,\% at 50\,Gy, 1\,\% at 54\,Gy and 10\,\% at 61\,Gy in conventional
fractionation \cite{kirkpatrick2010quantec}. The discrepancy is an order of
magnitude, it does not touch the equivalent dose, and it is a property of the
tabulated Lyman parameters rather than of the calculation. The same failure has
been reported clinically: the Lyman-Kutcher-Burman model with historical
parameters does not predict spinal cord tolerance in radiosurgery
\cite{daly2012lkb}.

\subsection*{What this establishes}

The equivalent doses are internally coherent and externally plausible: agreement
to 0.06\,Gy with HYPO-RT-PC, the correct ordering of the two CHHiP arms, a
linear tail that is conservative rather than permissive for PACE-B, and a
stereotactic cord calculation consistent with HyTEC. Two limits are equally
clear. Imposing a single prostate $\alpha/\beta$ of 3.1\,Gy understates the
equivalence PROFIT demonstrated, by about 5\,Gy. And the conventional spinal cord
complication curve is incompatible with QUANTEC by an order of magnitude.

No arithmetic defect is identified in the equivalent doses. The fixed
radiobiological library, however, cannot be presented as clinically validated:
the weaker point is the calibration of the complication parameters, and after it
the use of one prostate $\alpha/\beta$ for every case. Both are recorded in
the limitations section of the article and are why the complication and control
probabilities are reported as indicative.

\bibliographystyle{elsarticle-num}
\bibliography{refs}

\end{document}
